\documentclass[11pt,svgnames,smaller]{beamer}

\usepackage{CommandsAndStyle}

\usepackage{amsthm}
\usepackage{amsmath}
\usepackage{subfig}
\usepackage{xcolor}
\usepackage{textpos}
\usepackage{transparent}
\usepackage{tikz}
\usepackage{epigraph}
\usepackage[ruled,vlined,linesnumbered]{algorithm2e}
\usepackage{amsmath,amsfonts,amssymb,amsthm}
\usepackage{latexsym}
\usepackage{graphicx} % [demo] is just for the example
\usepackage{wrapfig}
\usepackage{array, makecell} %
\usepackage{hyperref}

\SetKwComment{Comment}{$\triangleright$\ }{}

\usepackage{caption}

\mode<presentation>

\author{Gaspare Ferraro}

\institute[University of Pisa]{University of Pisa\\Department of Computer Science}

\setbeamertemplate{section page}
{
	\begin{centering}
		\begin{beamercolorbox}[sep=16pt,center]{part title}
			\usebeamerfont{section title}\insertsection\par
		\end{beamercolorbox}
	\end{centering}
}

\title{TwoPaCo}
\subtitle{An efficient algorithm to build the compacted de Bruijn graph\\from many complete genomes}
\date{Pisa, 29 May 2018}
\titlegraphic{\vfill}
%\includegraphics[height=0.4cm]{images/creative_commons.png}}

\setbeamercolor{title}{fg=black!65!black}

\renewcommand{\partname}{Part}

\begin{document}
	
 \logo{\transparent{0.2}\includegraphics[height=3cm]{images/cherubino}}

	\begin{frame} 
	\titlepage
	\end{frame}
			
	\part{Introduction}
\section{Introduction}

%%%%%%%%%%%%%%%%%%%%%%%%%%%%%%%%%%%%%%%%%%%%%%%%%%%%%%%%%%%%%%%%%%%%%%%%%%%%%%%%%%%%%%%%%%%%%%%%%%%%%%%%%%%%%%%%%%%%%%%%%%%
%%%%%%%%%%%%%%%%%%%%%%%%%%%%%%%%%%%%%%%%%%%%%%%%%%%%%%%%%%%%%%%%%%%%%%%%%%%%%%%%%%%%%%%%%%%%%%%%%%%%%%%%%%%%%%%%%%%%%%%%%%%

\begin{frame}
	\partpage
	\centering
\end{frame}

%%%%%%%%%%%%%%%%%%%%%%%%%%%%%%%%%%%%%%%%%%%%%%%%%%%%%%%%%%%%%%%%%%%%%%%%%%%%%%%%%%%%%%%%%%%%%%%%%%%%%%%%%%%%%%%%%%%%%%%%%%%
%%%%%%%%%%%%%%%%%%%%%%%%%%%%%%%%%%%%%%%%%%%%%%%%%%%%%%%%%%%%%%%%%%%%%%%%%%%%%%%%%%%%%%%%%%%%%%%%%%%%%%%%%%%%%%%%%%%%%%%%%%%

\begin{frame}
	\frametitle{A pan-genomic algorithm}
	\centering
	
	From the greek word \textit{pan} (everything)
	
	\medskip
	
	Given a series of complete genomes
	
	\medskip

	\includegraphics[height=0.4cm]{images/gen1}

  \medskip
  	
	\includegraphics[height=0.4cm]{images/gen2}
	
	\pause
	
	\medskip

  What we want to see:
  
  \begin{itemize}
    \item {\color{orange}Similarities}
    \item {\color{purple}Differences}
    \item {\color{green}Duplications}
  \end{itemize}	
	
	\medskip
	
	\includegraphics[height=1cm]{images/gen3}
	
\end{frame}

%%%%%%%%%%%%%%%%%%%%%%%%%%%%%%%%%%%%%%%%%%%%%%%%%%%%%%%%%%%%%%%%%%%%%%%%%%%%%%%%%%%%%%%%%%%%%%%%%%%%%%%%%%%%%%%%%%%%%%%%%%%
%%%%%%%%%%%%%%%%%%%%%%%%%%%%%%%%%%%%%%%%%%%%%%%%%%%%%%%%%%%%%%%%%%%%%%%%%%%%%%%%%%%%%%%%%%%%%%%%%%%%%%%%%%%%%%%%%%%%%%%%%%%

\begin{frame}
	\frametitle{de Bruijn graph}
	\centering
	  
  Sequence \#1: {\color{red}TGACGTC} ($2$-mers: {\color{red} TG GA AC CG GT TC})
                            
  \medskip

	\includegraphics[height=0.7cm]{images/debruin-0}

  \medskip
  
  Sequence \#2: {\color{blue}TGACTTC} ($2$-mers: {\color{blue} TG GA AC CT TT TC})
	
	\medskip
	
	\includegraphics[height=0.7cm]{images/debruin-1}
	
	\pause

	\medskip
	
	Find common vertices

	\medskip
  \includegraphics[height=1.5cm]{images/debruin-2}

	\pause
	                           
  \medskip
  
  Merge them
	\medskip

	\includegraphics[height=1.3cm]{images/debruin-3}


\end{frame}

%%%%%%%%%%%%%%%%%%%%%%%%%%%%%%%%%%%%%%%%%%%%%%%%%%%%%%%%%%%%%%%%%%%%%%%%%%%%%%%%%%%%%%%%%%%%%%%%%%%%%%%%%%%%%%%%%%%%%%%%%%%
%%%%%%%%%%%%%%%%%%%%%%%%%%%%%%%%%%%%%%%%%%%%%%%%%%%%%%%%%%%%%%%%%%%%%%%%%%%%%%%%%%%%%%%%%%%%%%%%%%%%%%%%%%%%%%%%%%%%%%%%%%%

\begin{frame}
	\frametitle{Compacted de Bruijn graph}
	\centering
	
  To compact the de Bruijn graph we need to find the junctions
  
  \medskip
  
  A vertex is a junction if :
  
  Is the initial or final part of some sequence\\
  Has more than 1 in-edge or more than 1 out-edge\\
	
	\medskip

	\includegraphics[height=2cm]{images/debruin-3a}

  \pause

	\medskip

  Remove simple paths by connect junctions  	

	\medskip

	\includegraphics[height=1.5cm]{images/debruin-4}
	
\end{frame}

%%%%%%%%%%%%%%%%%%%%%%%%%%%%%%%%%%%%%%%%%%%%%%%%%%%%%%%%%%%%%%%%%%%%%%%%%%%%%%%%%%%%%%%%%%%%%%%%%%%%%%%%%%%%%%%%%%%%%%%%%%%
%%%%%%%%%%%%%%%%%%%%%%%%%%%%%%%%%%%%%%%%%%%%%%%%%%%%%%%%%%%%%%%%%%%%%%%%%%%%%%%%%%%%%%%%%%%%%%%%%%%%%%%%%%%%%%%%%%%%%%%%%%%

%\begin{frame}
%	\frametitle{Junctions}
%	\centering
%\end{frame}

%%%%%%%%%%%%%%%%%%%%%%%%%%%%%%%%%%%%%%%%%%%%%%%%%%%%%%%%%%%%%%%%%%%%%%%%%%%%%%%%%%%%%%%%%%%%%%%%%%%%%%%%%%%%%%%%%%%%%%%%%%%
%%%%%%%%%%%%%%%%%%%%%%%%%%%%%%%%%%%%%%%%%%%%%%%%%%%%%%%%%%%%%%%%%%%%%%%%%%%%%%%%%%%%%%%%%%%%%%%%%%%%%%%%%%%%%%%%%%%%%%%%%%%

\begin{frame}
	\frametitle{The problem}
	\centering
	
	How to construct the compacted de Bruijn graph,\\ from many complete genomes,\\ avoiding the ordinary graph construction?
	
	\pause
	
	\medskip
	
	$\downarrow$
	
	\medskip

  
  \textbf{Which $k$-mers are junctions?}
  
  \medskip
  
  +
  
  \medskip
  
  How to construct the compacted de Bruijn graph given the junctions?
  
	
\end{frame}

	\part{The algorithm}

%%%%%%%%%%%%%%%%%%%%%%%%%%%%%%%%%%%%%%%%%%%%%%%%%%%%%%%%%%%%%%%%%%%%%%%%%%%%%%%%%%%%%%%%%%%%%%%%%%%%%%%%%%%%%%%%%%%%%%%%%%%
%%%%%%%%%%%%%%%%%%%%%%%%%%%%%%%%%%%%%%%%%%%%%%%%%%%%%%%%%%%%%%%%%%%%%%%%%%%%%%%%%%%%%%%%%%%%%%%%%%%%%%%%%%%%%%%%%%%%%%%%%%%

\begin{frame}
	\partpage
	\centering
\end{frame}

%%%%%%%%%%%%%%%%%%%%%%%%%%%%%%%%%%%%%%%%%%%%%%%%%%%%%%%%%%%%%%%%%%%%%%%%%%%%%%%%%%%%%%%%%%%%%%%%%%%%%%%%%%%%%%%%%%%%%%%%%%%
%%%%%%%%%%%%%%%%%%%%%%%%%%%%%%%%%%%%%%%%%%%%%%%%%%%%%%%%%%%%%%%%%%%%%%%%%%%%%%%%%%%%%%%%%%%%%%%%%%%%%%%%%%%%%%%%%%%%%%%%%%%

\begin{frame}
	\frametitle{Naive algorithm}
	\centering
	
	\begin{itemize}
	  \item Store all $(k+1)$-mers in a hash table
	  \item For each $k$-mers query the possible edge
	  \item If only 1 in and 1 out edge, unmark as a junction
	\end{itemize}
	   
	\medskip
	
   \scalebox{.55}{                        
   \begin{algorithm}[H]
      \small
       \DontPrintSemicolon
       \SetKwInOut{Input}{Input}
       \SetKwInOut{Output}{Output}
       \Input{$S = \{s_{1}, ..., s_{m}\}$ genoma sequences\\$k$ integer, size of $k$-mers\\ $E$ empty set data structure\\$C$ Candidate set of junctions (\textbf{naively all positions are marked})}
       \Output{A reduce candidate set of junctions $C$}
       
       \BlankLine
       \ForEach{$s \in S$}{
        \For{$1 \leq i < |s| - k $}{
          \If{$C[s,i] = marked$}{ 
            $E \gets E \cup  s[i..i+k] \cup  s[i-1..i+k-1]$ \Comment*[r]{Store all $(k+1)$-mers}
          }
        }
       }
       \BlankLine
       \ForEach{$s \in S$}{
        \For{$1 \leq i < |s| - k $}{
          \If{$C[s,i] = marked$}{
          
            $(in, out) \gets (0, 0)$ \Comment*[r]{Count in/out edges}
            
            \ForEach{$c \in \{A, C, G, T\}$}{
              \If{$v \cdot c \in E$}{
                $in \gets in + 1$\;
              }
              \If{$c \cdot v \in E$}{
                $out \gets out + 1$\;
              }
            }
            
            \If{$(in, out) = (1,1)$}{
              $C[s,i] = unmarked$ \Comment*[r]{surely not a junction}
            }
          }
        }
       }
              
       \Return{$C$}
       \caption{\textsc{Filter-Junctions}}
   \end{algorithm}
   }
   
\end{frame}

%%%%%%%%%%%%%%%%%%%%%%%%%%%%%%%%%%%%%%%%%%%%%%%%%%%%%%%%%%%%%%%%%%%%%%%%%%%%%%%%%%%%%%%%%%%%%%%%%%%%%%%%%%%%%%%%%%%%%%%%%%%
%%%%%%%%%%%%%%%%%%%%%%%%%%%%%%%%%%%%%%%%%%%%%%%%%%%%%%%%%%%%%%%%%%%%%%%%%%%%%%%%%%%%%%%%%%%%%%%%%%%%%%%%%%%%%%%%%%%%%%%%%%%

\begin{frame}
	\frametitle{The memory issue}
	
	  First part of the naive algorithm:
	  
	  \bigskip
	  
	  \ForEach{$s \in S$}{
      \For{$1 \leq i < |s| - k $}{
          \If{$C[s,i] = marked$}{
            $E \gets E \cup  s[i..i+k] \cup  s[i-1..i+k-1]$\;
        }
      }
    }
     
	 \centering

    \bigskip

    We don't really need and, in almost all pratical cases, \\ we can't store all the possible $(k+1)$-mers.
    
    \bigskip
    
    Mainly because \textbf{only a little percentual} of them are \\ junction in the de Bruijn graph.
      
\end{frame}

%%%%%%%%%%%%%%%%%%%%%%%%%%%%%%%%%%%%%%%%%%%%%%%%%%%%%%%%%%%%%%%%%%%%%%%%%%%%%%%%%%%%%%%%%%%%%%%%%%%%%%%%%%%%%%%%%%%%%%%%%%%
%%%%%%%%%%%%%%%%%%%%%%%%%%%%%%%%%%%%%%%%%%%%%%%%%%%%%%%%%%%%%%%%%%%%%%%%%%%%%%%%%%%%%%%%%%%%%%%%%%%%%%%%%%%%%%%%%%%%%%%%%%%

\begin{frame}
	\frametitle{Bloom filter}
	
	\centering

	A space-efficient probabilistic hash table
	
	\medskip
	
	Bitmap $V$ of size $b$, $h$ hash functions $f_{0}, f_{1}, ..., f_{h-1} : U \rightarrow [0, b-1]$ \\
		
	\medskip
	
	insertion($x$) $\rightarrow V[f_{i}(x)] = 1$, $\forall$ $0 \leq i < h$ 
	
	\medskip
	
	contains($x$) $\rightarrow$ \textbf{probabily yes} if $V[f_{i}(x)] == 1$, $\forall$ $0 \leq i < h$
  \medskip
  
  \includegraphics[height=3cm]{images/bloom_filter}
  
  \medskip
	Probability of false positive, after $n$ insertion: $p_{FP} \simeq (1 - e^{-hn/m})^{h}$
	
	
\end{frame}

%%%%%%%%%%%%%%%%%%%%%%%%%%%%%%%%%%%%%%%%%%%%%%%%%%%%%%%%%%%%%%%%%%%%%%%%%%%%%%%%%%%%%%%%%%%%%%%%%%%%%%%%%%%%%%%%%%%%%%%%%%%
%%%%%%%%%%%%%%%%%%%%%%%%%%%%%%%%%%%%%%%%%%%%%%%%%%%%%%%%%%%%%%%%%%%%%%%%%%%%%%%%%%%%%%%%%%%%%%%%%%%%%%%%%%%%%%%%%%%%%%%%%%%

\begin{frame}
	\frametitle{Two Pass version}
	\centering
	
	\begin{itemize}
	  \item First pass: Select a set of junction candidates by insert all the $(k+1)$-mers in a bloom filter of choosing size
	  \item Second pass: Filter out the false positive by storing the reduce sets of $(k+1)$-mers in an hash table
	\end{itemize}

   \scalebox{.8}{                        
   \begin{algorithm}[H]
      \small
       \DontPrintSemicolon
       \SetKwInOut{Input}{Input}
       \SetKwInOut{Output}{Output}
       \Input{strings $S = \{s_{1}, ..., s_{m}\}$ genoma sequences\\integer $k$, size of $k$-mers\\integer $b$, size of bloom filter\\ Candidate set of junctions $C_{in}$ (\textbf{naively all positions are marked})}
       \Output{A reduce candidate set of junctions $C_{out}$}
       
       \BlankLine
       
       $F \gets$ empty bloom filter of size $b$ \Comment*[r]{First pass}
       $C_{temp} \gets \textsc{Filter-Junctions}(S, k, F, C_{in})$  \;

       \BlankLine

       $H \gets$ empty hash table \Comment*[r]{Second pass}
       $C_{out} \gets \textsc{Filter-Junctions}(S, k, H, C_{in})$ \;

       \BlankLine
              
       \Return{$C_{out}$}
       \caption{\textsc{Filter-Junctions-Two-Pass}}
   \end{algorithm}
   }
   

\end{frame}

%%%%%%%%%%%%%%%%%%%%%%%%%%%%%%%%%%%%%%%%%%%%%%%%%%%%%%%%%%%%%%%%%%%%%%%%%%%%%%%%%%%%%%%%%%%%%%%%%%%%%%%%%%%%%%%%%%%%%%%%%%%
%%%%%%%%%%%%%%%%%%%%%%%%%%%%%%%%%%%%%%%%%%%%%%%%%%%%%%%%%%%%%%%%%%%%%%%%%%%%%%%%%%%%%%%%%%%%%%%%%%%%%%%%%%%%%%%%%%%%%%%%%%%

\begin{frame}
	\frametitle{The memory issue$^{2}$}
	\centering
	
	How much memory do we use now?
	
	\begin{itemize}
	  \item First pass: 
	  \item Second pass: 
	\end{itemize}
	
	What if the hash table in the second pass 

\end{frame}

%%%%%%%%%%%%%%%%%%%%%%%%%%%%%%%%%%%%%%%%%%%%%%%%%%%%%%%%%%%%%%%%%%%%%%%%%%%%%%%%%%%%%%%%%%%%%%%%%%%%%%%%%%%%%%%%%%%%%%%%%%%
%%%%%%%%%%%%%%%%%%%%%%%%%%%%%%%%%%%%%%%%%%%%%%%%%%%%%%%%%%%%%%%%%%%%%%%%%%%%%%%%%%%%%%%%%%%%%%%%%%%%%%%%%%%%%%%%%%%%%%%%%%%

\begin{frame}

	\frametitle{Multiple rounds: dealing with memory restrictions}
	\centering
	
	\begin{itemize}
	  \item Partitionate input $k$-mers in 
	  \item ciao
	\end{itemize}
	
	\scalebox{0.7}{                        
   \begin{algorithm}[H]
      \small
       \DontPrintSemicolon
       \SetKwInOut{Input}{Input}
       \SetKwInOut{Output}{Output}
       \Input{strings $S = \{s_{1}, ..., s_{m}\}$ genoma sequences\\integer $k$, size of $k$-mers\\integer $b$, size of bloom filter\\$l$ }
       \Output{$C_{out}$ all the junctions in the compacted de Bruijn graph}
       
       \BlankLine
       
       $F \gets$ empty bloom filter of size $b$ \Comment*[r]{First pass}
       $C_{temp} \gets \textsc{Filter-Junctions}(S, k, F, C_{in})$  \;

       \BlankLine

       $H \gets$ empty hash table \Comment*[r]{Second pass}
       $C_{out} \gets \textsc{Filter-Junctions}(S, k, H, C_{in})$ \;

       \BlankLine
              
       \Return{$C_{out}$}
       \caption{\textsc{Filter-Junctions-Two-Pass}}
   \end{algorithm}
   }
   


\end{frame}

%%%%%%%%%%%%%%%%%%%%%%%%%%%%%%%%%%%%%%%%%%%%%%%%%%%%%%%%%%%%%%%%%%%%%%%%%%%%%%%%%%%%%%%%%%%%%%%%%%%%%%%%%%%%%%%%%%%%%%%%%%%
%%%%%%%%%%%%%%%%%%%%%%%%%%%%%%%%%%%%%%%%%%%%%%%%%%%%%%%%%%%%%%%%%%%%%%%%%%%%%%%%%%%%%%%%%%%%%%%%%%%%%%%%%%%%%%%%%%%%%%%%%%%

\begin{frame}
	\frametitle{Parallelization scheme}
	\centering
\end{frame}

	\part{Results}

\begin{frame}
	\partpage
	\centering
\end{frame}

	
	\section{Conclusion}
	
	\begin{frame}
  \frametitle{References}
  
  [1] Minkin, I., Pham, S., Medvedev, P. (2016). 
  
  \textbf{TwoPaCo}: An efficient algorithm to build the compacted de Bruijn graph from many complete genomes	  

  \medskip
  
  [2] Minkin, I., Patel, A., Kolmogorov, M., Vyahhi, N., Pham, S. (2013). 
  
  \textbf{Sibelia}: a scalable and comprehensive synteny block generation tool for closely related microbial
genomes.

  \medskip
  
  [3] Marcus, S., Lee, H., Schatz, M. C. (2014). 
  
  \textbf{SplitMEM}: a graphical algorithm for pan-genome analysis with suffix skips.

  \medskip

  [4] Baier, U., Beller, T., Ohlebusch, E. (2015). 
  
  Graphical pan-genome analysis with compressed suffix trees and the \textbf{Burrows-Wheeler transform}.


	\end{frame}

	\begin{frame}
		\frametitle{Conclusion}
		\centering
		\Large
		Thanks for your attention!
		
		\bigskip
		
		Questions?
 	\end{frame}
  
\end{document}
